% 第19批：学习指导自测题8计算题8.3.1--8.3.5；印刷173--174／PDF181--182。
\subsection{自测题8：静电能、旋转带电半环与电磁感应计算}

\sourcetest{8.3.1}{173}{181}
\problemtext 两金属导体球相距很远，半径均为10cm，分别带至电势200V和400V。随后用细导线相连并达到等势，求静电场力在电荷重新分配过程中所做的功。
\workings 远离的球近似各为对地独立电容 $C=4\pi\varepsilon_0 R$。接线后总电荷守恒，同半径故最终等势
\[V_f=\frac{CV_1+CV_2}{2C}=\frac{200+400}{2}=300\,{\rm V}.\]
原电场能为 $W_i=\tfrac12 C(V_1^2+V_2^2)$，终电场能为 $W_f=2(\tfrac12CV_f^2)=CV_f^2$。差额即电场释放的能量：
\[
\begin{aligned}
W_{\rm field}&=W_i-W_f=\frac C4(V_2-V_1)^2\\
&=\pi\varepsilon_0 R(200\,{\rm V})^2
=\boxed{1.11\times10^{-7}\ {\rm J}},
\end{aligned}
\]
取 $\varepsilon_0=8.85\times10^{-12}$SI。
\crosscheck 两球半径相同但初始电势不同，最终各为300V；真实连接电路中场能可以转化为焦耳热和电磁辐射，不应声称静电总能量守恒不变。

\sourcetest{8.3.2}{173}{181}
\problemtext 半径 $a$、线电荷密度为常数 $\lambda>0$ 的均匀带电细半圆环，以角速度 $\omega$ 绕它所在平面内、经过半圆两端点的直径轴匀速转动。求半圆圆心的磁感应强度以及整体磁矩大小。
\workings 设半圆上微元的极角为 $\theta\in[0,\pi]$，轴到微元距离 $\rho=a\sin\theta$，轴向偏移 $z=a\cos\theta$，电量 $dq=\lambda a\,d\theta$。旋转微元形成频率 $\omega/(2\pi)$ 的小电流环
\[dI=\frac{\omega\,dq}{2\pi}.\]
圆心磁场的轴向贡献为
\[dB=\frac{\mu_0dI\,\rho^2}{2(\rho^2+z^2)^{3/2}}
=\frac{\mu_0\lambda\omega}{4\pi}\sin^2\theta\,d\theta.\]
利用 $\int_0^\pi\sin^2\theta\,d\theta=\pi/2$，
\[\boxed{B_O=\frac{\mu_0\lambda\omega}{8}}.\]
每一小电流环的磁矩 $dP_m=dI\pi\rho^2=(\omega/2)\,dq\,\rho^2$，故
\[
\boxed{P_m=\frac{\omega\lambda a^3}{2}
\int_0^\pi\sin^2\theta\,d\theta
=\frac{\pi\lambda\omega a^3}{4}}.
\]
两向量对正电荷均沿角速度轴向。
\crosscheck 这不是将整根半圆弧当成一个固定平面圆电流环；它绕\emph{自身直径}转动，各微元扫过的圆半径不同，需按 $\sin^2\theta$ 积分。

\sourcetest{8.3.3}{174}{182}
\problemtext 空间存在平行同向的均匀电场 $\bm E$ 与磁场 $\bm B$；电子电量 $-e$、质量 $m$，初速 $v$ 与 $\bm E$ 夹角 $\alpha$。求加速度大小，并描述运动轨迹。
\workings 取 $\hat{\bm z}$ 沿 $\bm E,\bm B$，洛伦兹力为
\[m\dot{\bm v}=-e(\bm E+\bm v\times\bm B).\]
磁场只作用在垂直分量，保持 $v_\perp=v\sin\alpha$ 的大小；电场只改变平行速度，$v_z(t)=v\cos\alpha-(eE/m)t$。因两力互相正交，
\[
\boxed{|\bm a|=\frac e m\sqrt{E^2+B^2v^2\sin^2\alpha}},
\]
大小与时间无关。垂直方向以角速度 $\omega_c=eB/m$ 作圆周运动，半径 $\rho=mv\sin\alpha/(eB)$；轴向位置
\[\boxed{z(t)=v\cos\alpha\,t-\frac{eE}{2m}t^2}.\]
因此轨迹是\textbf{沿磁力线方向受匀加速运动调制的螺旋线}，螺距随时间变化。
\crosscheck 若 $E=0$ 为等螺距匀速螺旋；若 $\alpha=0$ 为沿场的直线变速运动。电子带负电，轴向电加速度与 $\bm E$ 方向相反。

\sourcetest{8.3.4}{174}{182}
\problemtext 无限长直导线沿竖直向上载恒定电流 $I$；其右侧同平面有直角三角形线圈 $ABC$：$BC=a$ 与导线垂直，$AC=b$ 与导线平行，最近顶点 $B$ 离导线 $d$。三角形线圈以水平速度 $v$ 向右远离导线。求总感应电动势大小及感应电流方向。
\workings 取距导线水平坐标 $x$，三角形竖向截面长度 $h(x)=b(x-d)/a$，$d\le x\le d+a$。直导线场为 $B(x)=\mu_0I/(2\pi x)$，垂直纸面向内。磁通绝对值
\[
\Phi(d)=\frac{\mu_0Ib}{2\pi a}
\int_d^{d+a}\frac{x-d}{x}\,dx
=\frac{\mu_0Ib}{2\pi a}
\left[a-d\ln\left(1+\frac ad\right)\right].
\]
远离过程中 $\dot d=v$，于是
\[
\boxed{|\mathcal E|=
\frac{\mu_0Ibv}{2\pi a}
\left[\ln\left(1+\frac ad\right)-\frac a{d+a}\right]}.
\]
原向里磁通减少，所以感应电流产生向里的磁场，方向为\textbf{顺时针}（沿 $A\to C\to B\to A$）。
\crosscheck 对数项必须减去 $a/(a+d)$；在远场 $d\gg a$，方括号近似 $a^2/(2d^2)>0$，电动势按 $d^{-2}$ 衰减。

\sourcetest{8.3.5}{174}{182}
\problemtext 无限长直导线载恒定电流 $I$，一根长度 $b$、与该导线平行的短金属棒初始距导线 $a$。金属棒以速率 $v$ 沿垂直于导线与棒初始所在平面的方向平移。求时刻 $t\ge0$ 的杆两端动生电动势大小。
\workings 令直导线沿 $+\hat{\bm y}$，棒初始位于 $x=a,z=0$，随后移到 $z=vt$，棒上任一点到直导线的距离 $\rho=\sqrt{a^2+v^2t^2}$。毕奥--萨伐尔定律
\[\bm B=\frac{\mu_0I}{2\pi(a^2+z^2)}
\left(z\,\hat{\bm x}-a\,\hat{\bm z}\right).\]
移动速度 $\bm v=v\hat{\bm z}$，其动生力 $\bm v\times\bm B$ 沿棒的 $y$ 方向，模为 $\mu_0 I v|z|/[2\pi(a^2+z^2)]$。沿长度 $b$ 积分，代入 $z=vt$：
\[
\boxed{|\mathcal E(t)|=
\frac{\mu_0I b v^2t}{2\pi(a^2+v^2t^2)}}\qquad(t\ge0).
\]
\crosscheck $t=0$ 时棒的移动速度恰与该处磁场平行，所以电动势为零；随远离运动先增大后减小，在 $t=a/v$ 达到极大值。
